\appendix
\section{Source-specific corrections and proof-status revisions}
\label{app:audit}

This appendix gives the principal proposed changes. The accompanying
\src{CORRECTIONS.txt} contains a fuller register, including exact
source line ranges and smaller notation repairs. All locators refer
to the commit on the title page. A false identity, an unproved claim,
and a correct statement with an incomplete proof are different
findings; they are identified separately below.

\subsection{The inverse-color mixed report}

\noindent\textbf{Source:}
\path{reports/eisenstein-gaussian-mixed-cuberoot-doubles__abd2e7116aa9.tex}.

The abstract, the table at source lines 170--184, and especially
\src{thm:parallel} at lines 219--226 claim two new constants in
each mixed tower. This is an \textbf{established error} relative to
the stated base family. Corollary~\ref{cor:four} gives explicit
single-data reductions of all four named candidates, and
Theorem~\ref{thm:mixed} eliminates any additional inverse-color
directions at every positive integral weight.

Replace the theorem by the all-weight quotient statement in
Section~\ref{sec:mixed}; update the abstract, table interpretation,
and any later dimension counts that use those directions.
The report's general numerical-evidence caveat does not repair a
false displayed theorem.
The claimed infinite asymptotic tail at lines 104--109 also needs a
finite truncation with a proved remainder, and the requested-digit
guarantee for Wynn acceleration at lines 111--118 needs an actual
error theorem. These are separate from the exact mixed reduction.
The missing upstream evaluator prevents a specific diagnosis of the
original search failure.

\subsection{Foundational multiple-polylogarithm conventions}

\noindent\textbf{Source:}
\path{articles/multiple-polylogarithms-introduction.tex}.

\paragraph{Convergence (lines 231--235).}
The stated absolute-convergence theorem is false already for
\(\Li_1(\ii)=\sum_{n\ge1}\ii^n/n\): the original series converges,
but its series of absolute values diverges.
A safe replacement in the stated positive-index setting is:
for \(|z_j|\le1\) and \(\Repart s_j\ge1\), absolute convergence
holds if \(\Repart s_1>1\) or \(|z_1|<1\).
More generally, strict inequalities
\(|z_1\cdots z_j|<1\) for every prefix give an interior domain of
absolute convergence. Boundary convergence and permissible
rearrangements should be stated separately. For positive integral
multiple-zeta indices, ordinary convergence is equivalent to
\(s_1\ge2\).

\paragraph{Iterated-integral orientation (lines 278--287).}
Use a consistent outermost-first recursion
\[
G(a_1,\ldots,a_w;z)=\int_0^z\frac{dt}{t-a_1}
                         G(a_2,\ldots,a_w;t),\qquad G(;z)=1.
\]
Its simplex is \(0<t_w<\cdots<t_1<1\), with \(a_j\) attached to
\(t_j\), and differentiation strips the \emph{first} letter.
The corresponding dictionary is
\begin{align*}
\Li_{s_1,\ldots,s_d}(z_1,\ldots,z_d)
=(-1)^dG(&0^{s_1-1},z_1^{-1},
          0^{s_2-1},(z_1z_2)^{-1},\ldots,\\
         &0^{s_d-1},(z_1\cdots z_d)^{-1};1).
\end{align*}
The arguments are inverse prefix products. The printed suffix
products are incompatible with this convention.
The same simplex reversal is needed in
\path{articles/gaussian-eisenstein-double-polylogs.tex},
\src{eq:G-def}--\src{eq:goncharov-dbl}, lines 141--146:
the printed orientation would make its proposed \(G(0,z^{-1};1)\)
representation diverge at the inner integration endpoint.

\paragraph{HPL recursion and signs (lines 491--555).}
A leading zero does not generally cause divergence:
\(H(0,1;z)=\Li_2(z)\) is a direct counterexample.
Use the recursion whenever its tail makes the integral convergent,
and define all-zero words separately by \(\log^w(z)/w!\).
With the source's convention
\(\zeta(\overline{s})=\sum_{n\ge1}(-1)^n/n^s\),
\[
H(0,-1;1)=-\Li_2(-1)=-\zeta(\overline2)=\frac{\pi^2}{12}.
\]
The sign in its identification with the barred zeta value must change.

\paragraph{Alphabet closure (lines 559--579).}
The full alphabet \(\{0,1,-1,\infty\}\) is not preserved by
every listed substitution. The weight-one counterexamples
\[
H(-1;1-z)=\log(2-z),\qquad
H(-1;z^2)=\log(1+z^2)
\]
introduce the singularities \(2\) and \(\pm\ii\), respectively.
The transformations \(-z\), \(1/z\), and
\((1-z)/(1+z)\) do preserve the singular set, with branches tracked.
Reflection \(1-z\) preserves the smaller \(\{0,1\}\) alphabet.
An algorithm acting on the full class must allow the necessary
enlarged alphabet.

\paragraph{Shuffle versus stuffle (lines 759--763).}
The correct identities are
\[
\Li_1(z)^2=2\Li_{1,1}(z,1)
          =2\Li_{1,1}(z,z)+\Li_2(z^2).
\]
The displayed formula with \(2\Li_{1,1}(z,z)\) alone is false:
that double starts at order \(z^3\), while the square starts at
order \(z^2\). The neighboring word example for
\(\zeta(2)\zeta(3)\) should shuffle the words \(01\) and \(001\),
of total weight five, rather than the printed weight-eight pair.

\paragraph{Unsupported numerical dimensions.}
The weight-five multiple-zeta table, lines 471--478, supplies
spanning reductions but not the rational independence of its two
generators. Replace the actual dimension assertion by an upper bound
or a stated conjectural value. Related sign and kernel repairs are
listed in the correction register; they do not invalidate every
low-weight identity in the introduction.

\subsection{Depth, parity, and proposed arithmetic bases}

\noindent\textbf{Sources:}
\path{articles/gaussian-multiple-polylog-depth.tex} and
\path{articles/gaussian-eisenstein-double-polylogs.tex}.

The definition of \(S_p\) in the first article, lines 128--132,
must require \(p\ge1\) for the ordinary series. At \(p=0\), its
terms fail to tend to zero.
Replace the \emph{precisely when} content of \src{thm:alternation},
lines 185--192, by the proved positive inclusion for every
\(S_{2m+1}\). Section~\ref{sec:euler} supplies its exact formula.
The source's \(S_5\) and \(S_7\) identities survive this change.
Its even-index nonmembership assertion is unproved; this was
already recognized in the intake README.

The uncolored analogy at line 351 includes
\(\zeta(1,2,1)\) and \(\zeta(1,1,2)\), which diverge as ordinary
positive-term multiple-zeta series. A regularized statement needs
an explicit regularization and a derivation.
Statements describing even beta values as new transcendental numbers
or a Champernowne canary as jointly algebraically independent are
not established by the supplied arguments.

In the second article, the dimension interpretations around
\src{thm:gauss-5} and \src{thm:eis-1} are \textbf{unsupported
arithmetic lower bounds}, not identities disproved here.
Replace them by the demonstrated spanning claims, or by an exact
formal corank if the complete matrix and certificate are supplied.
Products of depth-one and depth-two values lie in \(D^3\) in general:
\[
D^aD^b\subseteq D^{a+b}.
\]
The algebra generated by all depth-at-most-two values is a different
object from the linear filtration step \(D^2\).

The binomial depth-dimension formula near \src{thm:deligne},
lines 345--352, does not follow from free-Lie generation alone.
Word length in chosen Lie generators is not automatically MPL depth.
Even a positive ambient motivic dimension does not identify the
images of the three selected periods as nonzero or independent.
Adding a period conjecture cannot supply these missing motivic
calculations. Retain the proved reductions and reopen the specific
depth-three classification.

\subsection{Gamma laws, formal ranks, and complexity}

\noindent\textbf{Source:}
\path{articles/gamma-lattice-and-certificates.tex}.

The arithmetic completeness attributed to Koblitz--Ogus in
\src{law:ko}, lines 133--137, is not an available theorem of that
scope. The sufficient algebraicity criterion and the conjectural
converse must be distinguished; the intake README already flags
this issue. A primary discussion of the distinction is
\cite[\S1.4.2]{ABP}.

The unconditional replacement is:
\emph{The exact reducer is sound and complete for membership in the
rational span of the encoded reflection and multiplication rows.
Its free columns form a basis of that formal quotient.
Completeness for all multiplicative relations among actual gamma
values is not asserted.}
Section~\ref{sec:distribution} gives the formal anchored reflected
rank \(\varphi(N)/2\) for \(N\ge3\); at \(N=120\), this is 16.
The explicit identities accompanied by valid law certificates remain
valid.

The raw all-divisor row construction, lines 105--121, emits
\[
\left\lfloor\frac{N-1}{2}\right\rfloor
 +\sum_{\substack{d\mid N\\d\ge2}}\frac Nd
=\left\lfloor\frac{N-1}{2}\right\rfloor+\sigma_1(N)-N.
\]
This is not uniformly \(O(N)\), since \(\sigma_1(N)/N\) is
unbounded. State the exact count or a valid general upper bound;
row dependency does not retroactively change the raw output count.

\subsection{Stieltjes ranks and the complex-character correction}

\noindent\textbf{Sources:}
\path{articles/stieltjes-parameter-derivative-tower.tex},
\path{articles/stieltjes-antiderivative-ladder.tex}, and
\path{reports/stieltjes-derivative-relations__41afbe98f668.md}.

The derivative article's \src{thm:rank}, lines 195--212, states
an all-\(q,k\) conclusion with finite evidence.
The antiderivative article's \src{prop:jetrank}, lines 946--975,
more carefully states its finite range and conjectures the extension.
The missing all-denominator proof is now supplied in
Section~\ref{sec:distribution}. Replace the finite-experiment
justification by that proof, retaining the formal-quotient
qualification. The derivative article's gap was already an
intake warning.

The supporting report, section 4, lines 146 and 148, omits complex
conjugation. Put \(c_q=\gamma+\log(2\pi/q)-1\). The corrected
odd-character identity is
\[
\frac{L'(2,\chi)}{L(2,\chi)}
+\frac12\overline{\frac{L''(-1,\chi)}{L'(-1,\chi)}}=c_q.
\]
For an even primitive character the corrected second logarithmic
identity is
\[
(\log L)''(2,\chi)
-\overline{(\log L)''(-1,\chi)}=1+\frac{\pi^2}{12}.
\]
These are \textbf{established formula errors} for complex
characters, not merely stylistic omissions.
For the odd quartic character modulo five with \(\chi(2)=\ii\),
the uncorrected first formula has a discrepancy approximately
\(-0.1438652751124278\,\ii\).
For the even cubic character modulo seven with
\(\chi(3)=\omega\), the uncorrected second formula has discrepancy
\(-0.0339162597722613\,\ii\).
Section~\ref{sec:bridges} proves the corrected identities at every
order. The corresponding first-level bridges in the two main
Stieltjes articles already have the correct conjugation and should
not be altered on this account.

The claim of mutually independent second-derivative atoms in the
derivative article, lines 303--315, is not established by the PSLQ
experiments. Describe the retained constants as formal generators
unless an arithmetic lower-bound proof is supplied.
The antiderivative article should call its Hurwitz-zeta construction
\emph{a canonical analytic interpolant}, rather than claim uniqueness
from integer data: adding \(c\sin(2\pi z)\) preserves both integer
values and a unit-shift recurrence. Its other small corrections are
listed separately in the register.

\subsection{Polygamma parity and CM evidence}

\noindent\textbf{Sources:}
\path{articles/polylog-polygamma-bridge.tex} and
\path{articles/eisenstein-row-sums-cm-polygamma.tex}.

The Clausen component of the rational grid in the first article,
lines 149--151, has reflection parity \((-1)^{n+1}\) at
polylogarithm weight \(n\). It is not always the odd component.
Thus \(\Cl_2\) is represented by a trigamma difference, while
\(\Cl_3\) uses a tetragamma sum. The displayed laws already have
the correct alternating behavior; the prose needs repair.
The claimed exact converse classification of new Clausen atoms is
not established by the numerical searches. Its positive reductions
and exact \(L\)-function bridges remain useful.

In the CM article, \src{neg:sums} and \src{res:genus} use failed
searches to suggest irrationality or exact algebraic degree.
These are numerical observations. Even an exact quartic equation
only establishes degree at most four until irreducibility and root
identification are supplied.
The weight-two question is resolved in Section~\ref{sec:modular}.
The isolated imaginary trigamma value remains separate from the
modular aggregate identities.

The source's Hurwitz recurrence and its weight-twelve coefficients
are correct in the unnormalized convention:
\[
G_{12}=\frac{18G_4^3+25G_6^2}{143}.
\]
They should not be replaced by normalized Eisenstein coefficients
without carrying the zeta factors.

\subsection{A replacement rule for numerical claims}

Where the only evidence is a finite integer-relation search, use
the statement:
\emph{No relation was found for the specified basis, precision, and
search parameters. The listed constants are retained as formal
generators in this reduction scheme; arithmetic independence is not
established.}

A successful search proposes an identity until an exact proof or
certificate is supplied. A failed search is not automatically a
rigorous bounded-height exclusion either: that needs certified input
enclosures and a certified exclusion argument. A zero canary
coefficient does not prevent a false relation among the other
entries. These distinctions preserve the value of numerical
exploration while preventing unsupported theorem statements.

The older correction
\[
6\Li_2(1/3)-\Li_2(1/9)=\frac{\pi^2}{3}-\log^2 3
\]
to \path{reports/report-1.tex} was already recorded in the
repository README. It is not a new finding of this continuation.
Other pre-existing README caveats remain in force.
The present review does not certify every statement left unchanged.
